% Options for packages loaded elsewhere
\PassOptionsToPackage{unicode}{hyperref}
\PassOptionsToPackage{hyphens}{url}
\documentclass[
]{article}
\usepackage{xcolor}
\usepackage[margin=1in]{geometry}
\usepackage{amsmath,amssymb}
\setcounter{secnumdepth}{-\maxdimen} % remove section numbering
\usepackage{iftex}
\ifPDFTeX
  \usepackage[T1]{fontenc}
  \usepackage[utf8]{inputenc}
  \usepackage{textcomp} % provide euro and other symbols
\else % if luatex or xetex
  \usepackage{unicode-math} % this also loads fontspec
  \defaultfontfeatures{Scale=MatchLowercase}
  \defaultfontfeatures[\rmfamily]{Ligatures=TeX,Scale=1}
\fi
\usepackage{lmodern}
\ifPDFTeX\else
  % xetex/luatex font selection
\fi
% Use upquote if available, for straight quotes in verbatim environments
\IfFileExists{upquote.sty}{\usepackage{upquote}}{}
\IfFileExists{microtype.sty}{% use microtype if available
  \usepackage[]{microtype}
  \UseMicrotypeSet[protrusion]{basicmath} % disable protrusion for tt fonts
}{}
\makeatletter
\@ifundefined{KOMAClassName}{% if non-KOMA class
  \IfFileExists{parskip.sty}{%
    \usepackage{parskip}
  }{% else
    \setlength{\parindent}{0pt}
    \setlength{\parskip}{6pt plus 2pt minus 1pt}}
}{% if KOMA class
  \KOMAoptions{parskip=half}}
\makeatother
\usepackage{color}
\usepackage{fancyvrb}
\newcommand{\VerbBar}{|}
\newcommand{\VERB}{\Verb[commandchars=\\\{\}]}
\DefineVerbatimEnvironment{Highlighting}{Verbatim}{commandchars=\\\{\}}
% Add ',fontsize=\small' for more characters per line
\usepackage{framed}
\definecolor{shadecolor}{RGB}{248,248,248}
\newenvironment{Shaded}{\begin{snugshade}}{\end{snugshade}}
\newcommand{\AlertTok}[1]{\textcolor[rgb]{0.94,0.16,0.16}{#1}}
\newcommand{\AnnotationTok}[1]{\textcolor[rgb]{0.56,0.35,0.01}{\textbf{\textit{#1}}}}
\newcommand{\AttributeTok}[1]{\textcolor[rgb]{0.13,0.29,0.53}{#1}}
\newcommand{\BaseNTok}[1]{\textcolor[rgb]{0.00,0.00,0.81}{#1}}
\newcommand{\BuiltInTok}[1]{#1}
\newcommand{\CharTok}[1]{\textcolor[rgb]{0.31,0.60,0.02}{#1}}
\newcommand{\CommentTok}[1]{\textcolor[rgb]{0.56,0.35,0.01}{\textit{#1}}}
\newcommand{\CommentVarTok}[1]{\textcolor[rgb]{0.56,0.35,0.01}{\textbf{\textit{#1}}}}
\newcommand{\ConstantTok}[1]{\textcolor[rgb]{0.56,0.35,0.01}{#1}}
\newcommand{\ControlFlowTok}[1]{\textcolor[rgb]{0.13,0.29,0.53}{\textbf{#1}}}
\newcommand{\DataTypeTok}[1]{\textcolor[rgb]{0.13,0.29,0.53}{#1}}
\newcommand{\DecValTok}[1]{\textcolor[rgb]{0.00,0.00,0.81}{#1}}
\newcommand{\DocumentationTok}[1]{\textcolor[rgb]{0.56,0.35,0.01}{\textbf{\textit{#1}}}}
\newcommand{\ErrorTok}[1]{\textcolor[rgb]{0.64,0.00,0.00}{\textbf{#1}}}
\newcommand{\ExtensionTok}[1]{#1}
\newcommand{\FloatTok}[1]{\textcolor[rgb]{0.00,0.00,0.81}{#1}}
\newcommand{\FunctionTok}[1]{\textcolor[rgb]{0.13,0.29,0.53}{\textbf{#1}}}
\newcommand{\ImportTok}[1]{#1}
\newcommand{\InformationTok}[1]{\textcolor[rgb]{0.56,0.35,0.01}{\textbf{\textit{#1}}}}
\newcommand{\KeywordTok}[1]{\textcolor[rgb]{0.13,0.29,0.53}{\textbf{#1}}}
\newcommand{\NormalTok}[1]{#1}
\newcommand{\OperatorTok}[1]{\textcolor[rgb]{0.81,0.36,0.00}{\textbf{#1}}}
\newcommand{\OtherTok}[1]{\textcolor[rgb]{0.56,0.35,0.01}{#1}}
\newcommand{\PreprocessorTok}[1]{\textcolor[rgb]{0.56,0.35,0.01}{\textit{#1}}}
\newcommand{\RegionMarkerTok}[1]{#1}
\newcommand{\SpecialCharTok}[1]{\textcolor[rgb]{0.81,0.36,0.00}{\textbf{#1}}}
\newcommand{\SpecialStringTok}[1]{\textcolor[rgb]{0.31,0.60,0.02}{#1}}
\newcommand{\StringTok}[1]{\textcolor[rgb]{0.31,0.60,0.02}{#1}}
\newcommand{\VariableTok}[1]{\textcolor[rgb]{0.00,0.00,0.00}{#1}}
\newcommand{\VerbatimStringTok}[1]{\textcolor[rgb]{0.31,0.60,0.02}{#1}}
\newcommand{\WarningTok}[1]{\textcolor[rgb]{0.56,0.35,0.01}{\textbf{\textit{#1}}}}
\usepackage{graphicx}
\makeatletter
\newsavebox\pandoc@box
\newcommand*\pandocbounded[1]{% scales image to fit in text height/width
  \sbox\pandoc@box{#1}%
  \Gscale@div\@tempa{\textheight}{\dimexpr\ht\pandoc@box+\dp\pandoc@box\relax}%
  \Gscale@div\@tempb{\linewidth}{\wd\pandoc@box}%
  \ifdim\@tempb\p@<\@tempa\p@\let\@tempa\@tempb\fi% select the smaller of both
  \ifdim\@tempa\p@<\p@\scalebox{\@tempa}{\usebox\pandoc@box}%
  \else\usebox{\pandoc@box}%
  \fi%
}
% Set default figure placement to htbp
\def\fps@figure{htbp}
\makeatother
\setlength{\emergencystretch}{3em} % prevent overfull lines
\providecommand{\tightlist}{%
  \setlength{\itemsep}{0pt}\setlength{\parskip}{0pt}}
\usepackage{bookmark}
\IfFileExists{xurl.sty}{\usepackage{xurl}}{} % add URL line breaks if available
\urlstyle{same}
\hypersetup{
  pdftitle={Linear mixed models with lme4},
  hidelinks,
  pdfcreator={LaTeX via pandoc}}

\title{Linear mixed models with lme4}
\author{}
\date{\vspace{-2.5em}}

\begin{document}
\maketitle

\textbf{Inference lab} using the five-step template: Hypothesis →
Visualise → Assumptions → Conduct → Conclude.

\subsection{Learning objectives}\label{learning-objectives}

\begin{itemize}
\tightlist
\item
  Fit random-intercept and random-slope linear mixed models to
  longitudinal data.
\item
  Test whether adding a random slope improves fit (likelihood ratio test
  with ML).
\item
  Interpret REML estimates of variance components.
\end{itemize}

\subsection{Prerequisites}\label{prerequisites}

Linear regression; familiarity with \texttt{lm()}.

\subsection{Background}\label{background}

Longitudinal data have two sources of variation: differences between
people and differences over time within a person. A random-intercept
model lets each person have their own average level; a random-slope
model lets each person also have their own rate of change. When the
slopes genuinely vary, omitting the random slope underestimates the
standard error of the fixed-effect slope and produces \emph{p}-values
that are too small.

REML (restricted maximum likelihood) is the default because it gives
unbiased estimates of the variance components by integrating out the
fixed effects. It is the right choice for reporting variance components.
For comparing models that differ in fixed effects, you must refit with
ML (likelihood ratio tests on REML fits are invalid); for comparing
models that differ only in random effects, REML is fine.

A common misstep is to interpret the variance of the random intercept as
the ``within-subject variance''. It is the between- subject variance of
the intercept. The within-subject variance is the residual variance, σ²,
once the random effects are accounted for.

\subsection{Setup}\label{setup}

\begin{Shaded}
\begin{Highlighting}[]
\FunctionTok{library}\NormalTok{(tidyverse)}
\FunctionTok{library}\NormalTok{(lme4)}
\FunctionTok{library}\NormalTok{(lmerTest)}
\FunctionTok{set.seed}\NormalTok{(}\DecValTok{42}\NormalTok{)}
\FunctionTok{theme\_set}\NormalTok{(}\FunctionTok{theme\_minimal}\NormalTok{(}\AttributeTok{base\_size =} \DecValTok{12}\NormalTok{))}
\end{Highlighting}
\end{Shaded}

\subsection{1. Hypothesis}\label{hypothesis}

In the \texttt{sleepstudy} data, reaction time increases with days of
sleep deprivation, and the \emph{rate} of increase varies by subject.

\subsection{2. Visualise}\label{visualise}

\begin{Shaded}
\begin{Highlighting}[]

\NormalTok{sleepstudy }\SpecialCharTok{|\textgreater{}}
  \FunctionTok{ggplot}\NormalTok{(}\FunctionTok{aes}\NormalTok{(Days, Reaction, }\AttributeTok{group =}\NormalTok{ Subject)) }\SpecialCharTok{+}
  \FunctionTok{geom\_line}\NormalTok{(}\AttributeTok{alpha =} \FloatTok{0.3}\NormalTok{) }\SpecialCharTok{+}
  \FunctionTok{geom\_smooth}\NormalTok{(}\FunctionTok{aes}\NormalTok{(}\AttributeTok{group =} \DecValTok{1}\NormalTok{), }\AttributeTok{method =} \StringTok{"lm"}\NormalTok{, }\AttributeTok{se =} \ConstantTok{FALSE}\NormalTok{,}
              \AttributeTok{colour =} \StringTok{"firebrick"}\NormalTok{) }\SpecialCharTok{+}
  \FunctionTok{labs}\NormalTok{(}\AttributeTok{x =} \StringTok{"Days of deprivation"}\NormalTok{, }\AttributeTok{y =} \StringTok{"Reaction time (ms)"}\NormalTok{)}
\end{Highlighting}
\end{Shaded}

\subsection{3. Assumptions}\label{assumptions}

Linear relationship between Days and Reaction for each subject; normally
distributed random intercepts and slopes (if included); residuals
independent within and across subjects after random effects account for
clustering; approximately constant residual variance.

\subsection{4. Conduct}\label{conduct}

\begin{Shaded}
\begin{Highlighting}[]
\NormalTok{           data }\OtherTok{=}\NormalTok{ sleepstudy, REML }\OtherTok{=} \ConstantTok{TRUE}\ErrorTok{)}
\NormalTok{rs }\OtherTok{\textless{}{-}} \FunctionTok{lmer}\NormalTok{(Reaction }\SpecialCharTok{\textasciitilde{}}\NormalTok{ Days }\SpecialCharTok{+}\NormalTok{ (Days }\SpecialCharTok{|}\NormalTok{ Subject),}
           \AttributeTok{data =}\NormalTok{ sleepstudy, }\AttributeTok{REML =} \ConstantTok{TRUE}\NormalTok{)}
\FunctionTok{summary}\NormalTok{(rs)}
\end{Highlighting}
\end{Shaded}

\begin{Shaded}
\begin{Highlighting}[]
\CommentTok{\# but we use REFIT = TRUE in anova() to get correct comparison.}
\FunctionTok{anova}\NormalTok{(ri, rs, }\AttributeTok{refit =} \ConstantTok{FALSE}\NormalTok{)}
\end{Highlighting}
\end{Shaded}

\begin{Shaded}
\begin{Highlighting}[]
\NormalTok{fm\_ml }\OtherTok{\textless{}{-}} \FunctionTok{lmer}\NormalTok{(Reaction }\SpecialCharTok{\textasciitilde{}}\NormalTok{ Days }\SpecialCharTok{+}\NormalTok{ (Days }\SpecialCharTok{|}\NormalTok{ Subject),}
              \AttributeTok{data =}\NormalTok{ sleepstudy, }\AttributeTok{REML =} \ConstantTok{FALSE}\NormalTok{)}
\NormalTok{fm\_ml}
\end{Highlighting}
\end{Shaded}

\begin{Shaded}
\begin{Highlighting}[]
\NormalTok{sleepstudy }\SpecialCharTok{|\textgreater{}}
  \FunctionTok{mutate}\NormalTok{(}\AttributeTok{fit\_rs =} \FunctionTok{predict}\NormalTok{(rs)) }\SpecialCharTok{|\textgreater{}}
  \FunctionTok{ggplot}\NormalTok{(}\FunctionTok{aes}\NormalTok{(Days, Reaction, }\AttributeTok{group =}\NormalTok{ Subject)) }\SpecialCharTok{+}
  \FunctionTok{geom\_point}\NormalTok{(}\AttributeTok{alpha =} \FloatTok{0.4}\NormalTok{) }\SpecialCharTok{+}
  \FunctionTok{geom\_line}\NormalTok{(}\FunctionTok{aes}\NormalTok{(}\AttributeTok{y =}\NormalTok{ fit\_rs), }\AttributeTok{colour =} \StringTok{"steelblue"}\NormalTok{) }\SpecialCharTok{+}
  \FunctionTok{facet\_wrap}\NormalTok{(}\SpecialCharTok{\textasciitilde{}}\NormalTok{ Subject, }\AttributeTok{ncol =} \DecValTok{6}\NormalTok{) }\SpecialCharTok{+}
  \FunctionTok{theme}\NormalTok{(}\AttributeTok{strip.text =} \FunctionTok{element\_text}\NormalTok{(}\AttributeTok{size =} \DecValTok{7}\NormalTok{))}
\end{Highlighting}
\end{Shaded}

\subsection{5. Concluding statement}\label{concluding-statement}

\begin{quote}
In the sleepstudy data, reaction time increased by about
\texttt{round(fixef(rs){[}"Days"{]},\ 1)} ms per day of deprivation,
with substantial subject-to-subject variation in slope (SD ≈
\texttt{round(sqrt(VarCorr(rs)\$Subject{[}2,\ 2{]}),\ 1)} ms/day).
Adding a random slope significantly improved fit over a random-intercept
model (likelihood-ratio test \emph{p} \textless{} 0.001).
\end{quote}

\subsection{Common pitfalls}\label{common-pitfalls}

\begin{itemize}
\tightlist
\item
  Testing random-effects with REML-based LRTs without refitting.
\item
  Ignoring the random slope when slopes clearly vary.
\item
  Treating the \texttt{Intercept} variance as the within-subject error.
\item
  Using \emph{p}-values from the Wald \emph{z} on small samples.
\end{itemize}

\subsection{Further reading}\label{further-reading}

\begin{itemize}
\tightlist
\item
  Bates D et al.~(2015), \emph{Fitting linear mixed-effects models using
  lme4}.
\item
  Pinheiro JC, Bates DM. \emph{Mixed-Effects Models in S and S-PLUS}.
\item
  Gelman A, Hill J. \emph{Data Analysis Using Regression and
  Multilevel/Hierarchical Models}.
\end{itemize}

\subsection{Session info}\label{session-info}

\begin{Shaded}
\begin{Highlighting}[]

\end{Highlighting}
\end{Shaded}

\subsection{See also --- chapter index}\label{see-also-chapter-index}

\begin{itemize}
\tightlist
\item
  \href{../index.Rmd}{Methods in this chapter}
\item
  \href{../../../ch00-find-your-method.Rmd}{Find your method (Chapter
  0)}
\end{itemize}

\end{document}
